% preamble.tex
\usepackage{luatexja}
\usepackage{luatexja-ruby}
\usepackage{luatexja-fontspec}

% === Fuentes japonesas ===
\setmainjfont{Harano Aji Mincho}
\setsansjfont{Harano Aji Gothic}

% === Emoji fix: Emoji package ===
% Usaremos el paquete emoji directamente para todos los emojis.
\usepackage{emoji}
\setemojifont{Segoe UI Emoji}

% === Colores ===
\usepackage{xcolor}
\definecolor{jpred}{HTML}{BC002D}
\definecolor{jpblue}{HTML}{1A3A6B}
\definecolor{jpgold}{HTML}{C8A951}
\definecolor{sakura}{HTML}{FFB7C5}
\definecolor{lightgray}{HTML}{F7F7F7}
\definecolor{darkgray}{HTML}{555555}
\definecolor{softgreen}{HTML}{D4EDDA}

% === Diseño de página ===
\usepackage[a4paper, top=2.5cm, bottom=2.5cm, left=3cm, right=3cm]{geometry}

% === Tablas ===
\usepackage{booktabs}
\usepackage{tabularx}
\usepackage{array}
\usepackage{longtable}

% === tcolorbox ===
\usepackage[most]{tcolorbox}
\tcbuselibrary{skins, breakable}

\tcbset{
  vocabbox/.style={
    enhanced, breakable,
    colback=lightgray, colframe=jpblue,
    coltitle=white, fonttitle=\bfseries\large,
    attach boxed title to top left={yshift=-2mm, xshift=5mm},
    boxed title style={colback=jpblue, rounded corners},
    rounded corners
  },
  grammarbox/.style={
    enhanced, breakable,
    colback=white, colframe=jpred,
    coltitle=white, fonttitle=\bfseries\large,
    attach boxed title to top left={yshift=-2mm, xshift=5mm},
    boxed title style={colback=jpred, rounded corners},
    rounded corners
  },
  phrasebox/.style={
    enhanced,
    colback=sakura!15, colframe=sakura!80!black,
    rounded corners, fonttitle=\bfseries
  },
  kanjibox/.style={
    enhanced, breakable,
    colback=jpgold!8, colframe=jpgold!70!black,
    coltitle=white, fonttitle=\bfseries\large,
    attach boxed title to top left={yshift=-2mm, xshift=5mm},
    boxed title style={colback=jpgold!70!black, rounded corners},
    rounded corners
  },
  ejerciciobox/.style={
    enhanced,
    colback=softgreen!50, colframe=softgreen!80!black,
    rounded corners, fonttitle=\bfseries,
    coltitle=darkgray
  }
}

% === Hyperref ===
\usepackage{hyperref}
\hypersetup{
  colorlinks=true,
  linkcolor=jpblue,
  urlcolor=jpred,
  pdftitle={Apuntes de Japonés N5-N4},
  pdfauthor={Aldo ZM}
}

% === Títulos ===
\usepackage{titlesec}

\titleformat{\chapter}[display]
  {\normalfont\huge\bfseries\color{jpred}}{}
  {0pt}{\huge}
  [\vspace{0.3em}\color{jpred}\hrule height 1.5pt\vspace{0.3em}]

\titleformat{\section}
  {\large\bfseries\color{jpblue}}{}
  {0em}{}
  [\vspace{0.1em}\color{jpblue!40}\hrule height 0.4pt]

\titleformat{\subsection}
  {\normalsize\bfseries\color{darkgray}}{}
  {0em}{}

% =============================================
% COMANDOS CUSTOM
% =============================================

\newcommand{\unidadheader}[4]{%
  \begin{tcolorbox}[
    enhanced, colback=jpred!5, colframe=jpred,
    coltitle=white, fonttitle=\bfseries\Large,
    title={\emoji{crossed-flags} Unidad #1 \enspace \ruby{#2}{} \enspace #3},
    attach boxed title to top center={yshift=-3mm},
    boxed title style={colback=jpred, rounded corners}
  ]
  \textit{\textcolor{jpblue}{\emoji{speaking-head} Tema comunicativo: #4}}
  \end{tcolorbox}
  \vspace{0.5em}
}

\newcommand{\vocabitem}[4]{%
  \ruby{#1}{#2} & #3 & \textcolor{darkgray}{\small\textit{#4}} \\[4pt]
}

\newcommand{\patron}[3]{%
  \begin{tcolorbox}[grammarbox, title={\emoji{gear} #1}]
  \begin{center}{\Large\textbf{#2}}\end{center}
  \vspace{0.3em}
  #3
  \end{tcolorbox}
}

\newcommand{\ej}[2]{%
  \par\vspace{0.2em}
  \hspace{1em}▶ #1 \hfill {\small\textit{#2}}
  \vspace{0.1em}
}

\newcommand{\frase}[2]{%
  \begin{tcolorbox}[phrasebox]
  \textbf{#1}\\[2pt]
  {\small\textit{#2}}
  \end{tcolorbox}
  \vspace{0.2em}
}

\newcommand{\kanjirow}[5]{%
  {\LARGE #1} & #2 & #3 & #4 & #5 \\[4pt]
}

% Reemplazos manuales de caracteres unicode crudos en el texto a comandos del paquete emoji
\usepackage{newunicodechar}
\newunicodechar{🎌}{\emoji{crossed-flags}}
\newunicodechar{🗣}{\emoji{speaking-head}}
\newunicodechar{📌}{\emoji{pushpin}}
\newunicodechar{⚙}{\emoji{gear}}
\newunicodechar{📝}{\emoji{memo}}
\newunicodechar{🀄}{\emoji{mahjong-red-dragon}}
\newunicodechar{🔢}{\emoji{input-numbers}}
\newunicodechar{👆}{\emoji{backhand-index-pointing-up}}
\newunicodechar{🗺}{\emoji{world-map}}
\newunicodechar{⏰}{\emoji{alarm-clock}}
\newunicodechar{❤}{\emoji{red-heart}}
\newunicodechar{🍱}{\emoji{bento-box}}
\newunicodechar{🛍}{\emoji{shopping-bags}}
\newunicodechar{🎨}{\emoji{artist-palette}}
\newunicodechar{👨}{\emoji{man}}
\newunicodechar{👩}{\emoji{woman}}
\newunicodechar{👧}{\emoji{girl}}
\newunicodechar{👦}{\emoji{boy}}
\newunicodechar{⌛}{\emoji{hourglass-done}}
\newunicodechar{☀️}{\emoji{sun}}
\newunicodechar{☀️}{\emoji{sun}} % fix possible variation
\newunicodechar{☀}{\emoji{sun}}
\newunicodechar{🔗}{\emoji{link}}
\newunicodechar{🔄}{\emoji{counterclockwise-arrows-button}}
\newunicodechar{🌟}{\emoji{glowing-star}}
\newunicodechar{📅}{\emoji{calendar}}
\newunicodechar{⚖}{\emoji{balance-scale}}
\newunicodechar{🚆}{\emoji{train}}
\newunicodechar{🏥}{\emoji{hospital}}
\newunicodechar{💪}{\emoji{flexed-biceps}}
\newunicodechar{🌠}{\emoji{shooting-star}}
\newunicodechar{🎁}{\emoji{wrapped-gift}}
\newunicodechar{💬}{\emoji{speech-balloon}}
\newunicodechar{💭}{\emoji{thought-balloon}}
\newunicodechar{🔀}{\emoji{shuffle-tracks-button}}
\newunicodechar{📢}{\emoji{loudspeaker}}
\newunicodechar{🙇}{\emoji{person-bowing}}
\newunicodechar{📑}{\emoji{bookmark-tabs}}
\newunicodechar{📈}{\emoji{chart-increasing}}
\newunicodechar{🏁}{\emoji{chequered-flag}}
\newunicodechar{📊}{\emoji{bar-chart}}
\newunicodechar{📇}{\emoji{card-index}}
\newunicodechar{🇯}{\emoji{Japan}} % J
\newunicodechar{🇵}{\emoji{Japan}} % P  (J+P together usually make the flag, but we'll map them individually just in case to the crossed flags if needed, or leave them. Better to just map the JP flag if it's a single unicode char).
\newunicodechar{✗}{\textcolor{jpred}{X}}
\newunicodechar{📚}{\emoji{books}}
